\documentclass[11pt]{article}
\usepackage[utf8]{inputenc}
\usepackage{geometry}
\geometry{a4paper, margin=1in}
\usepackage{hyperref}
\usepackage{booktabs}
\usepackage{enumitem}
\usepackage{amsmath}
\usepackage{float}
\usepackage{xcolor}

\title{Lifecycle 2 --- Q\&A Preparation \\ \large Questions Your Guide Will Ask}
\author{Study Reference}
\date{2026-03-18}

\begin{document}

\maketitle

\tableofcontents
\newpage

% ============================================================================
\section{Temperature Scaling Questions}
% ============================================================================

\subsection{Q: What is temperature scaling? Why T = 1.1511?}

\textbf{Problem:} EfficientNet-B3 was overconfident --- saying 95\% when only 75\% accurate.

\textbf{Solution:} Divide logits by temperature $T$ before softmax:
\[
\hat{q}_i = \frac{\exp(z_i / T)}{\sum_j \exp(z_j / T)}
\]

\begin{itemize}[noitemsep]
    \item $T > 1$: Softens distribution (reduces overconfidence) --- our case
    \item $T < 1$: Sharpens distribution (increases confidence)
    \item $T = 1$: No change
\end{itemize}

\textbf{How T was found:} Optimized on validation set by minimizing negative log-likelihood using L-BFGS optimizer. Optimal value: $T = 1.1511$.

\textbf{Why it doesn't change accuracy:} Dividing ALL logits by same $T$ doesn't change which is largest. The argmax (predicted class) stays the same. Only probability magnitudes change.

\subsection{Q: What is calibration? Why does it matter?}

A model is ``calibrated'' when its confidence matches reality. If it says 80\% confidence on 100 images, approximately 80 should be correctly classified.

\textbf{Why it matters for our project:}
\begin{itemize}[noitemsep]
    \item Without calibration: ``95\% landslide'' might be wrong 25\% of the time --- dangerous
    \item With calibration: ``80\% landslide'' means it's right $\sim$80\% of the time --- actionable
    \item Enables meaningful thresholds: ``auto-alert if $>$80\%, human review if 50--80\%''
\end{itemize}

\subsection{Q: If temperature scaling doesn't change metrics, why bother?}

Because \textbf{trust matters}. In disaster monitoring:
\begin{enumerate}[noitemsep]
    \item Response teams need to know HOW confident the model really is
    \item Enables tiered response: high confidence $\to$ immediate action, medium $\to$ human verification
    \item Makes the ensemble fusion more reliable --- calibrated probabilities combine better than overconfident ones
\end{enumerate}

% ============================================================================
\section{Ensemble Questions}
% ============================================================================

\subsection{Q: Why ensemble? Why 60/40 and not 50/50?}

\textbf{Why ensemble:} Two models make different mistakes. Averaging cancels random errors. When EfficientNet is confused, AlexNet might be correct, and vice versa.

\textbf{Why 60/40:}
\begin{itemize}[noitemsep]
    \item EfficientNet-B3 is stronger (higher F1, higher AUC)
    \item 60\% weight lets it dominate when both agree
    \item AlexNet's 40\% can still override when EfficientNet is wrong
    \item 50/50 would give too much weight to the weaker model
\end{itemize}

\textbf{If asked ``why not learn the weights?'':} Valid future improvement. A meta-learner on validation predictions could optimize the ratio. For now, 60/40 heuristic worked well.

\subsection{Q: How does ensemble prediction work exactly?}

Step by step:
\begin{enumerate}[noitemsep]
    \item Same image goes through both models independently
    \item EfficientNet outputs softmax probabilities (after temperature scaling): e.g., [0.30, 0.70]
    \item AlexNet outputs softmax probabilities: e.g., [0.45, 0.55]
    \item Weighted average: $0.6 \times [0.30, 0.70] + 0.4 \times [0.45, 0.55] = [0.36, 0.64]$
    \item If landslide probability (0.64) $\geq$ threshold (0.467) $\to$ LANDSLIDE DETECTED
\end{enumerate}

\subsection{Q: Why is 0.467 the optimal threshold and not 0.5?}

0.5 is arbitrary. The optimal threshold maximizes the F1-score of the landslide class:

\begin{itemize}[noitemsep]
    \item We scan thresholds from 0.0 to 1.0 in small steps
    \item At each threshold, compute precision, recall, and F1 for the landslide class
    \item 0.467 gives the best F1 trade-off: more landslides caught (higher recall) at acceptable false alarm rate
    \item Lower than 0.5 because in disaster monitoring, we prefer to catch more landslides (even at cost of some false alarms)
\end{itemize}

% ============================================================================
\section{Vision Transformer Questions}
% ============================================================================

\subsection{Q: What is a Vision Transformer (ViT)?}

\textbf{CNN:} Slides local filters across image. Sees small patches. Builds global view gradually through many layers.

\textbf{ViT:} Splits image into 16$\times$16 patches, flattens them into vectors, processes with self-attention. Every patch ``attends to'' every other patch from layer 1.

\begin{table}[H]
    \centering
    \renewcommand{\arraystretch}{1.3}
    \begin{tabular}{lll}
        \toprule
        & \textbf{CNN} & \textbf{ViT} \\
        \midrule
        Receptive field & Local (grows with depth) & Global (from layer 1) \\
        Good at & Textures, edges, local patterns & Spatial relationships, context \\
        Weakness & Misses global context & May miss fine textures \\
        Input & 300$\times$300 pixels & 224$\times$224 $\to$ 14$\times$14 patches \\
        \bottomrule
    \end{tabular}
\end{table}

\subsection{Q: What is self-attention?}

Self-attention lets each patch look at every other patch and decide which are important:

\begin{enumerate}[noitemsep]
    \item Each patch gets three vectors: Query ($Q$), Key ($K$), Value ($V$)
    \item Attention score: $\text{Attention}(Q, K, V) = \text{softmax}\left(\frac{QK^T}{\sqrt{d_k}}\right) V$
    \item High score = ``these patches are related''
\end{enumerate}

\textbf{Example:} A debris patch at the hill bottom attends to a scar patch at the top --- the model understands they're part of the same landslide, even though far apart.

\subsection{Q: Why did ViT give a bigger improvement than AlexNet?}

Because CNN + CNN share the same blind spots, while CNN + ViT have \textbf{different inductive biases}:

\begin{itemize}[noitemsep]
    \item r4 (EfficientNet + AlexNet): Both CNNs. When one fails on global context, the other likely fails too. F1 = 72.97\%
    \item r5 (EfficientNet + ViT): Different architectures. CNN catches textures, ViT catches spatial layout. F1 = 75.98\% (+3.01\%)
\end{itemize}

\textbf{Key lesson:} \textit{Diversity of architecture} matters more than \textit{number of models}.

\subsection{Q: How did you train ViT-B/16?}

\begin{itemize}[noitemsep]
    \item Pre-trained on ImageNet (1.2M images)
    \item Replaced classification head: 1000 classes $\to$ 6 classes (our dataset)
    \item Froze all base parameters, only trained the new head initially
    \item Fine-tuned for 15 epochs on our training set
    \item Val accuracy: 83.50\% (standalone, before ensemble)
    \item Input: 224$\times$224 (ViT's native resolution)
\end{itemize}

% ============================================================================
\section{HMM Upgrade Questions}
% ============================================================================

\subsection{Q: Why increase hidden states from 4 to 8?}

4 states captured broad categories (Shallow, Deep, Debris, Rockfall). But landslides are triggered by different mechanisms:
\begin{itemize}[noitemsep]
    \item A monsoon-triggered mudslide has different recurrence patterns than an earthquake-triggered rockfall
    \item 8 states: Shallow Rain Slide, Deep Seismic Slide, Debris Flow, Rockfall, Mixed Flow, Complex Event, Monsoon Mudslide, Human-Triggered Slide
    \item More states = finer-grained temporal regimes = better forecasts
\end{itemize}

\subsection{Q: Why combine type $\times$ trigger into 42 symbols?}

Previously: HMM only observed landslide TYPE (7 symbols: Landslide, Mudslide, Rockfall, etc.). It didn't know WHY the landslide happened.

Now: Each observation is TYPE $\times$ TRIGGER (e.g., ``Mudslide + Heavy Rain'', ``Rockfall + Earthquake''). This gives the HMM environmental context.

$7 \text{ types} \times 6 \text{ triggers} = 42 \text{ combined symbols}$

\textbf{Result:} The HMM can now distinguish ``monsoon season debris flows'' from ``earthquake season rockfalls'' and make more targeted forecasts.

% ============================================================================
\section{Results Questions}
% ============================================================================

\subsection{Q: What was the biggest improvement in LC2?}

Replacing AlexNet with ViT-B/16 (r4 $\to$ r5):
\begin{itemize}[noitemsep]
    \item Accuracy: +2.15\% (82.40\% $\to$ 84.55\%)
    \item F1: +3.01\% (72.97\% $\to$ 75.98\%)
    \item AUC: +1.67\% (89.83\% $\to$ 91.50\%)
\end{itemize}

This was the single largest improvement in the entire project, confirming that architectural diversity in ensembles is the key.

\subsection{Q: Why did r3 have the same metrics as r2?}

Temperature scaling doesn't change predictions --- it only fixes confidence scores. The argmax (predicted class) is identical, so accuracy, precision, recall, F1, and AUC are unchanged. The improvement is in \textit{calibration quality}, not classification metrics.

\subsection{Q: What are the remaining weaknesses after LC2?}

Three gaps that motivate LC3:
\begin{enumerate}[noitemsep]
    \item \textbf{Class imbalance:} 72/28 split still biases training. 40 false negatives remain.
    \item \textbf{Single-pass fragility:} Same image, different orientation $\to$ different confidence. Some FNs were at 44--46\% (just below 46.7\% threshold).
    \item \textbf{Black box:} No way to verify WHY the model made a prediction.
\end{enumerate}

\end{document}
